\documentclass[a5j]{ltjarticle}

% For “『』” usage, see
% https://ja.wikipedia.org/wiki/%E6%8B%AC%E5%BC%A7#%E5%90%8D%E7%A7%B0%E3%81%AE%E6%8F%90%E7%A4%BA.
\newcommand*{\titleContent}{『\ruby{基|礎}{き|そ}日本語１』の\ruby{文|字|起|こし}{も|じ|お|}}
\newcommand*{\authorContent}{趙磊}

\title{\titleContent}
\author{\authorContent}

\usepackage{luatexja-ruby}

\usepackage{enumitem}
\newlist{dialogue}{description}{1}
\setlist[dialogue]{font=\normalfont,align=right,labelsep=.5\zw,labelwidth=2\zw,leftmargin=!,nosep}

% For “：” usage, see https://kotonoha-sha.jp/mojiokoshi/sample/ and
% https://satomasaki.com/column/manner-of-speaking/6809/.
\newcommand*{\YamadaSpeaks}{\item[山田：]}
\newcommand*{\ZhangSpeaks}{\item[張：]}
\newcommand*{\TanakaSpeaks}{\item[田中：]}

\usepackage[pagestyles]{titlesec}

\titleformat{\section}[hang]{\normalfont\Large\bfseries}{第\kansuji\value{section}課}{\zw}{}

\ltjsetparameter{kansujichar={0,`０}}
\ltjsetparameter{kansujichar={1,`１}}
\ltjsetparameter{kansujichar={2,`２}}
\ltjsetparameter{kansujichar={3,`３}}
\ltjsetparameter{kansujichar={4,`４}}
\ltjsetparameter{kansujichar={5,`５}}
\ltjsetparameter{kansujichar={6,`６}}
\ltjsetparameter{kansujichar={7,`７}}
\ltjsetparameter{kansujichar={8,`９}}
\ltjsetparameter{kansujichar={9,`９}}

\begin{document}
\maketitle

\section{どうぞよろしく}

\subsection*{本文}

\begin{dialogue}[labelwidth=\zw]
\item[A] 私は田中です。どうぞよろしく。
\item[B] 私は張です。どうぞよろしく。
\item
\item[A] これは何ですか。
\item[B] それはラジオです。
\item[A] それは何ですか。
\item[B] これはパソコンです。
\item[A] あれは何ですか。
\item[B] あれは雑誌です。
\item[A] あれは日本の雑誌ですか。
\item[B] はい、そうです。
\item[A] どれがラジカセですか。
\item[B] あれがラジカセです。\\
  あれは私のではありません。友達のです。
\item
\item[A] あの人は誰ですか。
\item[B] どの人ですか。
\item[A] あの女の人です。
\item[B] あの人は山田さんです。\\
  山田さんは東南大学の留学生です。\\
  山田さんは中国人ではありません。日本人です。
\item[A] 田中さんは誰ですか。
\item[B] 田中さんはこの人です。
\item[A] どなたが先生ですか。
\item[B] その男の人が先生です。
\end{dialogue}

\subsection*{会話（１）}

\begin{dialogue}
  \YamadaSpeaks
  はじめまして。
  \ZhangSpeaks
  はじめまして。
  \YamadaSpeaks
  山田です。
  \ZhangSpeaks
  張です。
  \YamadaSpeaks
  よろしくお願いします。
  \ZhangSpeaks
  こちらこそどうぞよろしくお願いします。
\end{dialogue}

\subsection*{会話（２）}

\begin{dialogue}
  \YamadaSpeaks
  張さん、こちらは田中さんです。\\
  私の友達です。
  \TanakaSpeaks
  田中です。どうぞよろしく。
  \YamadaSpeaks
  こちらは張です。\\
  張さんも私の友達です。
  \ZhangSpeaks
  張です。どうぞよろしく。\\
  田中さんは会社員ですか。
  \TanakaSpeaks
  いいえ、そうではありません。\\
  私は留学生です。東南大学の留学生です。\\
  張さんは？
  \ZhangSpeaks
  私も東南大学の学生です。\\
  日本語科の一年生です。どうぞよろしく。
\end{dialogue}

\section{}



\section{}



\section{}



\section{}



\section{}



\section{}



\section{}



\section{}



\section{}



\section{}



\section{}



\section{}



\section{}



\section{}



\section{}



\section{}



\section{}



\section{}



\section{}



\end{document}

% Local Variables:
% TeX-engine: luatex
% End:
